%
% Permission is granted to copy, distribute and/or modify this
% document under the terms of the Creative Common by-nc-sa License
% version 3.0 (CC BY-NC-SA 3.0). A copy of the license can be found at
% http://creativecommons.org/licenses/by-nc-sa/3.0/legalcode.
%

\usepackage[french]{babel}
\usepackage{multicol}

% Highlight macros
\newcommand{\highlight}[1]{\textcolor{structure.fg}{\bfseries #1}}

%% Title, subtitle, authors, institute, date, ...
\title{Analyse de malware statique et dynamique}
\subtitle{Les outils et quelques cas pratiques}

\author[E. Grelet \& A. Guémon]{Erwan Grelet \& Amélie Guémon\\[1em]}

%\author[E. Grelet \& A. Guémon]{Erwan Grelet \& Amélie Guémon\\[1em]
%  \texttt{\scriptsize <erwan.grelet@etu.u-bordeaux.fr>}\\[.3em]
%  \texttt{\scriptsize <amelie.guemon@etu.u-bordeaux.fr>}}

\institute[Master CSI, France]{Master CSI, Université de Bordeaux, France}

\date{\today}

%%%%%%%%%%%%%%%%%%%%%%%%%%[ Document ]%%%%%%%%%%%%%%%%%%%%%%%%%%
\begin{document}

\begin{frame}
  \vspace{3.5em}
  \titlepage
  \begin{center}
    \includegraphics[scale=.2]{cc-by-nc-sa.pdf}
  \end{center}
\end{frame}
%%%%%%%%%%%%%%%%%%%%%%%%%%%%%%%%%%%%%%%%%%%%%%%%%%%%%%%%%%%%%%%%%%%%%%
\section*{Introduction}

\begin{frame}[fragile]
  \frametitle{Introduction}
  \vspace{-.5em}
  Quelques statistiques \ldots
  \vfill
  \begin{itemize}\setlength\itemsep{1.5em}
        \item{58.31\% des emails échangés durant l'année 2016 étaient des spams,
        \\soit 3 points de plus qu'en 2015;}
        %https://securelist.com/analysis/kaspersky-security-bulletin/77483/kaspersky-security-bulletin-spam-and-phishing-in-2016/
        \item{Plus de 90\% des emails de phishing contenaient des ransomwares;}
        %sur le premier semestre 2016
        \item{En 2016, 1 088 900 victimes de banking trojans,
        30\% de plus qu'en 2015;}
        %https://securelist.com/analysis/publications/77623/financial-cyberthreats-in-2016/
        %\item{notamment pokémongo Trojan-Banker.AndroidOS.Tordow et guide ztorg
        \item{En septembre 2016, attaque DDOS de 1 Tb/s
        sur les serveurs d'OVH France;}
        %depuis IoT (152 464 caméras, routeurs, thermostats, \ldots
        \item{Code source du botnet Mirai disponible.}
        %avec scanner de vulnérabilité d'IoT intégré, \ldots
  \end{itemize}
  \vspace{-.25em}
\end{frame}
% ###################################################%
\begin{frame}[fragile]
  \frametitle{Introduction}
  \begin{multicols}{2}[{\large Quelques types de menaces:}]
  \begin{itemize}\setlength\itemsep{1.5em}
        \item{Adware;}
        \item{Spyware;}
        \item{Worm;}
        \item{Trojan;}
        \item{Rootkit;}
        \item{Backdoors;}
        \item{Keyloggers;}
        \item{Rogue security software;}
        \item{Ransomware;}
        \item{Browser Hijacker;}
        \item{Botnet;}
        \item{\ldots}
  \end{itemize}
  \end{multicols}
  \vspace{-.25em}
\end{frame}
% ###################################################%
\begin{frame}
  \frametitle{Plan}
  \tableofcontents[]
\end{frame}
%%%%%%%%%%%%%%%%%%%%%%%%%%%%%%%%%%%%%%%%%%%%%%%%%%%%%%%%%%%%%%%%%%%%%%
\section{Les Techniques d'Analyses}

\begin{frame}<handout:0>
  \frametitle{Plan}
  \tableofcontents[currentsection]
\end{frame}
%###################################################%
\begin{frame}[fragile]
  \frametitle{Les Techniques d'Analyses}
\vspace{-.25em}
2 techniques d'analyse: statique \& dynamique
\end{frame}
%###################################################%
\begin{frame}[fragile]
  \frametitle{Les Techniques d'Analyses -- Statique}
\vspace{-.25em}
Description analyse statique (Header + Strings)
\end{frame}
%###################################################%
\begin{frame}[fragile]
  \frametitle{Les Techniques d'Analyses -- Statique}
\vspace{-.25em}
Description analyse statique (Packer + Lecture de code)
\end{frame}
%###################################################%
\begin{frame}[fragile]
  \frametitle{Les Techniques d'Analyses -- Dynamique}
\vspace{-.25em}
Description analyse dynamique (Vm + INetSim + Analyse trafic)
\end{frame}
%###################################################%
\begin{frame}[fragile]
  \frametitle{Les Techniques d'Analyses -- Dynamique}
\vspace{-.25em}
Description analyse dynamique (Listing processus + Tracing + Débugueur)
\end{frame}
%%%%%%%%%%%%%%%%%%%%%%%%%%%%%%%%%%%%%%%%%%%%%%%%%%%%%%%%%%%%%%%%%%%%%%
\section{BillGates Botnet}

\begin{frame}<handout:0>
  \frametitle{Plan}
  \tableofcontents[currentsection]
\end{frame}
%###################################################%
\begin{frame}[fragile]
  \frametitle{BillGates Botnet -- Miscellaneous}
\vspace{-.25em}
Historique + infection pour IoT
\end{frame}
%###################################################%
\begin{frame}[fragile]
  \frametitle{BillGates Botnet -- Main}
\vspace{-.25em}
Main + cha\^ine obfusqu\'ee
\end{frame}
%###################################################%
\begin{frame}[fragile]
  \frametitle{BillGates Botnet -- Modules}
\vspace{-.25em}
Description Modules + sch\'ema
\end{frame}
%###################################################%
\begin{frame}[fragile]
  \frametitle{BillGates Botnet -- Attaques}
  \vspace{-.25em}
  Types d'attaques mises en place
\end{frame}
%%%%%%%%%%%%%%%%%%%%%%%%%%%%%%%%%%%%%%%%%%%%%%%%%%%%%%%%%%%%%%%%%%%%%%
\section{SageCrypt Ransomware}

\begin{frame}<handout:0>
  \frametitle{Plan}
  \tableofcontents[currentsection]
\end{frame}

%###################################################%
\begin{frame}[fragile]
  \frametitle{SageCrypt Ransomware -- Miscellaneous }
\vspace{-.25em}
Historique (CryLocker, Petya, \ldots)
\end{frame}
%###################################################%
\begin{frame}[fragile]
  \frametitle{SageCrypt Ransomware -- Miscellaneous }
\vspace{-.25em}
Infection par macro word + flood UAC + chiffre pas les fichiers important
+ layout
\end{frame}
%###################################################%
\begin{frame}[fragile]
  \frametitle{SageCrypt Ransomware -- Unpacking}
\vspace{-.25em}
Unpacking
\end{frame}
%###################################################%
\begin{frame}[fragile]
  \frametitle{SageCrypt Ransomware -- Chiffrement}
\vspace{-.25em}
Principe de chiffrement
\end{frame}
%%%%%%%%%%%%%%%%%%%%%%%%%%%%%%%%%%%%%%%%%%%%%%%%%%%%%%%%%%%%%%%%%%%%%%
\section*{Conclusion}

\begin{frame}[fragile]
  \frametitle{Conclusion}
\vspace{-.25em}

\end{frame}
%%%%%%%%%%%%%%%%%%%%%%%%%%%%%%%%%%%%%%%%%%%%%%%%%%%%%%%%%%%%%%%%%%%%%%
%\section{Références \& Lectures supplémentaires}

%\begin{frame}<handout:0>
%  \frametitle{Plan}
%  \tableofcontents[currentsection,subsectionstyle=hide]
%\end{frame}

%\nocite{*}
%\bibliographystyle{alpha}

%\begin{frame}[allowframebreaks]
%  \frametitle{Livres et références}
%  \bibliography{bibliography}
%\end{frame}

%%%%%%%%%%%%%%%%%%%%%%%%%%%%%%%%%%%%%%%%%%%%%%%%%%%%%%%%%%%%%%%%%%%%%%
\begin{frame}
  \vfill
  \centering
  \highlight{\Huge Questions~?}
  \vfill
\end{frame}
